\documentclass[10pt,a4paper]{amsart}
\usepackage{amsmath,amssymb,amsthm,graphicx}
\usepackage{verbatim}
\newenvironment{reference}{\comment}{\endcomment}
%\newenvironment{reference}{}{}
\newenvironment{slogan}{\comment}{\endcomment}
\newenvironment{history}{\comment}{\endcomment}

% For commutative diagrams we use Xy-pic
\usepackage[all]{xy}

% We use 2cell for 2-commutative diagrams.
\xyoption{2cell}
% Theorem environments.
%
\theoremstyle{plain}
\newtheorem{theorem}[subsection]{Theorem}
\newtheorem{proposition}[subsection]{Proposition}
\newtheorem{lemma}[subsection]{Lemma}

\theoremstyle{definition}
\newtheorem{definition}[subsection]{Definition}
\newtheorem{example}[subsection]{Example}
\newtheorem{exercise}[subsection]{Exercise}
\newtheorem{situation}[subsection]{Situation}

\theoremstyle{remark}
\newtheorem{remark}[subsection]{Remark}
\newtheorem{remarks}[subsection]{Remarks}

\numberwithin{equation}{subsection}

% Macros
%
\def\lim{\mathop{\mathrm{lim}}\nolimits}
\def\colim{\mathop{\mathrm{colim}}\nolimits}
\def\Spec{\mathop{\mathrm{Spec}}}
\def\Hom{\mathop{\mathrm{Hom}}\nolimits}
\def\Ext{\mathop{\mathrm{Ext}}\nolimits}
\def\SheafHom{\mathop{\mathcal{H}\!\mathit{om}}\nolimits}
\def\SheafExt{\mathop{\mathcal{E}\!\mathit{xt}}\nolimits}
\def\Sch{\mathit{Sch}}
\def\Mor{\mathop{\mathrm{Mor}}\nolimits}
\def\Ob{\mathop{\mathrm{Ob}}\nolimits}
\def\Sh{\mathop{\mathit{Sh}}\nolimits}
\def\NL{\mathop{N\!L}\nolimits}
\def\CH{\mathop{\mathrm{CH}}\nolimits}
\def\proetale{{pro\text{-}\acute{e}tale}}
\def\etale{{\acute{e}tale}}
\def\QCoh{\mathit{QCoh}}
\def\Ker{\mathop{\mathrm{Ker}}}
\def\Im{\mathop{\mathrm{Im}}}
\def\Coker{\mathop{\mathrm{Coker}}}
\def\Coim{\mathop{\mathrm{Coim}}}

% Boxtimes
%
\DeclareMathSymbol{\boxtimes}{\mathbin}{AMSa}{"02}

%
% Macros for moduli stacks/spaces
%
\def\QCohstack{\mathcal{QC}\!\mathit{oh}}
\def\Cohstack{\mathcal{C}\!\mathit{oh}}
\def\Spacesstack{\mathcal{S}\!\mathit{paces}}
\def\Quotfunctor{\mathrm{Quot}}
\def\Hilbfunctor{\mathrm{Hilb}}
\def\Curvesstack{\mathcal{C}\!\mathit{urves}}
\def\Polarizedstack{\mathcal{P}\!\mathit{olarized}}
\def\Complexesstack{\mathcal{C}\!\mathit{omplexes}}
% \Pic is the operator that assigns to X its picard group, usage \Pic(X)
% \Picardstack_{X/B} denotes the Picard stack of X over B
% \Picardfunctor_{X/B} denotes the Picard functor of X over B
\def\Pic{\mathop{\mathrm{Pic}}\nolimits}
\def\Picardstack{\mathcal{P}\!\mathit{ic}}
\def\Picardfunctor{\mathrm{Pic}}
\def\Deformationcategory{\mathcal{D}\!\mathit{ef}}
\usepackage{cleveref}
\crefname{originalreference}{original reference}{original references}
\Crefname{originalreference}{Original reference}{Original references}
\makeatletter
\newlabel{lemma-derived-quasi-coherent-small-etale-site}{{071Q}{source}{}{sourcepointer-1}{}}
\newlabel{lemma-derived-quasi-coherent-small-etale-site@cref}{{[originalreference][0][]071Q}{[1][0][]source}{}{}{}}
\makeatother
\crefname{originalreference}{original reference}{original references}
\Crefname{originalreference}{Original reference}{Original references}
\makeatletter
\newlabel{lemma-descend-perfect}{{08HG}{source}{}{sourcepointer-2}{}}
\newlabel{lemma-descend-perfect@cref}{{[originalreference][0][]08HG}{[1][0][]source}{}{}{}}
\makeatother
\crefname{originalreference}{original reference}{original references}
\Crefname{originalreference}{Original reference}{Original references}
\makeatletter
\newlabel{lemma-quasi-coherence-pushforward-direct-sums}{{08DZ}{source}{}{sourcepointer-3}{}}
\newlabel{lemma-quasi-coherence-pushforward-direct-sums@cref}{{[originalreference][0][]08DZ}{[1][0][]source}{}{}{}}
\makeatother
\crefname{originalproposition}{original proposition}{original propositions}
\Crefname{originalproposition}{Original proposition}{Original propositions}
\makeatletter
\newlabel{perfect-proposition-compact-is-perfect}{{09M1}{source}{}{sourcepointer-4}{}}
\newlabel{perfect-proposition-compact-is-perfect@cref}{{[originalproposition][0][]09M1}{[1][0][]source}{}{}{}}
\makeatother
\crefname{originallemma}{original lemma}{original lemmas}
\Crefname{originallemma}{Original lemma}{Original lemmas}
\makeatletter
\newlabel{sites-cohomology-lemma-dual-perfect-complex}{{08JJ}{source}{}{sourcepointer-5}{}}
\newlabel{sites-cohomology-lemma-dual-perfect-complex@cref}{{[originallemma][0][]08JJ}{[1][0][]source}{}{}{}}
\makeatother
\crefname{originallemma}{original lemma}{original lemmas}
\Crefname{originallemma}{Original lemma}{Original lemmas}
\makeatletter
\newlabel{sites-cohomology-lemma-perfect-independent-representative}{{08G6}{source}{}{sourcepointer-6}{}}
\newlabel{sites-cohomology-lemma-perfect-independent-representative@cref}{{[originallemma][0][]08G6}{[1][0][]source}{}{}{}}
\makeatother
\title{Compact objects of the derived quasi-coherent category are exactly the perfect ones}
\author{The Stacks Project Authors}
\date{stacks-project:spaces-perfect.tex:local-snapshot-20261012; GNU FDL 1.2 or later}
\begin{document}
\maketitle
\noindent\textit{Editorial note.} All delivered mathematical text is the author's own, copied from the pinned Stacks Project chapter file spaces-perfect.tex; no statement or proof step was rewritten, reordered or removed.\par\medskip
\noindent\textit{Editorial note.} The counted proof is the author's complete argument: the proof establishes compactness of perfect objects through the duality formula with direct-sum commutation, then reduces the converse to affine etale opens via Rj\_* direct-sum commutation and finally to the scheme case via the comparison lemmas\par\medskip
\noindent\textit{Editorial note.} The label proposition-compact-is-perfect names the counted statement, and the proof environment immediately following it in the source is delivered verbatim as its proof.\par\medskip
\noindent\textit{Editorial note.} The counted statement is carried by the author's own proposition environment in the chapter source.\par\medskip
\noindent\textit{Allowed standard prerequisites (editorial).}\begin{itemize}
\item the duality formula Hom\_\{D(O\_X)\}(K, M) = H\textasciicircum{}0(X, K\textasciicircum{}vee tensor\textasciicircum{}L M) for K perfect with dual K\textasciicircum{}vee (Cohomology on Sites, Lemma sites-cohomology-lemma-dual-perfect-complex)
\item commutation of tensoring with a perfect complex with direct sums, and commutation of H\textasciicircum{}0(X, -) with direct sums on D\_QCoh(O\_X) (lemma-quasi-coherence-pushforward-direct-sums of the same chapter)
\item perfectness of an object of D\_QCoh checked etale locally (Cohomology on Sites, Lemma sites-cohomology-lemma-perfect-independent-representative)
\item Rj\_* commutes with direct sums for quasi-compact separated morphisms j (lemma-quasi-coherence-pushforward-direct-sums of the same chapter) and the adjunction between restriction and Rj\_*
\item comparison of D\_QCoh on the small etale site with the scheme case (lemma-derived-quasi-coherent-small-etale-site and lemma-descend-perfect of the same chapter) and the scheme proposition perfect-proposition-compact-is-perfect of Derived Categories of Schemes
\end{itemize}\medskip
\ifdefined\hypertarget\hypertarget{sourcepointer-1}{}\fi
\noindent\textit{Source reference.} Original reference 071Q: sources/stacks-spaces-perfect/pinned-original-source.tex, source lines 402--402. Prerequisite pointer invoked at proof line(s) 4133; it names a standard fact of the same chapter that the counted argument uses, and every step of the counted proof built on it is retained verbatim in the delivered body.\par\medskip
\ifdefined\hypertarget\hypertarget{sourcepointer-2}{}\fi
\noindent\textit{Source reference.} Original reference 08HG: sources/stacks-spaces-perfect/pinned-original-source.tex, source lines 3288--3288. Prerequisite pointer invoked at proof line(s) 4134; it names a standard fact of the same chapter that the counted argument uses, and every step of the counted proof built on it is retained verbatim in the delivered body.\par\medskip
\ifdefined\hypertarget\hypertarget{sourcepointer-3}{}\fi
\noindent\textit{Source reference.} Original reference 08DZ: sources/stacks-spaces-perfect/pinned-original-source.tex, source lines 863--863. Prerequisite pointer invoked at proof line(s) 4115; it names a standard fact of the same chapter that the counted argument uses, and every step of the counted proof built on it is retained verbatim in the delivered body.\par\medskip
\ifdefined\hypertarget\hypertarget{sourcepointer-4}{}\fi
\noindent\textit{Source reference.} Original proposition 09M1: sources/stacks-chow/tags.txt, source lines 12164--12164. Prerequisite pointer invoked at proof line(s) 4137; it names a standard fact of a sibling chapter that the counted argument uses, and every step of the counted proof built on it is retained verbatim in the delivered body.\par\medskip
\ifdefined\hypertarget\hypertarget{sourcepointer-5}{}\fi
\noindent\textit{Source reference.} Original lemma 08JJ: sources/stacks-chow/tags.txt, source lines 10852--10852. Prerequisite pointer invoked at proof line(s) 4106; it names a standard fact of a sibling chapter that the counted argument uses, and every step of the counted proof built on it is retained verbatim in the delivered body.\par\medskip
\ifdefined\hypertarget\hypertarget{sourcepointer-6}{}\fi
\noindent\textit{Source reference.} Original lemma 08G6: sources/stacks-chow/tags.txt, source lines 10734--10734. Prerequisite pointer invoked at proof line(s) 4123; it names a standard fact of a sibling chapter that the counted argument uses, and every step of the counted proof built on it is retained verbatim in the delivered body.\par\medskip
\begin{proposition}
\label{proposition-compact-is-perfect}
Let $S$ be a scheme.
Let $X$ be a quasi-compact and quasi-separated algebraic space over $S$.
An object of $D_\QCoh(\mathcal{O}_X)$ is compact
if and only if it is perfect.
\end{proposition}

\begin{proof}
If $K$ is a perfect object of $D(\mathcal{O}_X)$ with dual
$K^\vee$ (Cohomology on Sites, Lemma
\ref{sites-cohomology-lemma-dual-perfect-complex})
we have
$$
\Hom_{D(\mathcal{O}_X)}(K, M) =
H^0(X, K^\vee \otimes_{\mathcal{O}_X}^\mathbf{L} M)
$$
functorially in $M$. Since $K^\vee \otimes_{\mathcal{O}_X}^\mathbf{L} -$
commutes with direct sums and since $H^0(X, -)$ commutes with direct
sums on $D_\QCoh(\mathcal{O}_X)$ by
Lemma \ref{lemma-quasi-coherence-pushforward-direct-sums}
we conclude that $K$ is compact in $D_\QCoh(\mathcal{O}_X)$.

\medskip\noindent
Conversely, let $K$ be a compact object of $D_\QCoh(\mathcal{O}_X)$.
To show that $K$ is perfect, it suffices to show that
$K|_U$ is perfect for every affine scheme $U$ \'etale over $X$, see
Cohomology on Sites, Lemma
\ref{sites-cohomology-lemma-perfect-independent-representative}.
Observe that $j : U \to X$ is a quasi-compact and separated morphism.
Hence
$Rj_* : D_\QCoh(\mathcal{O}_U) \to D_\QCoh(\mathcal{O}_X)$
commutes with direct sums, see
Lemma \ref{lemma-quasi-coherence-pushforward-direct-sums}.
Thus the adjointness of restriction to $U$ and $Rj_*$ implies that
$K|_U$ is a perfect object of $D_\QCoh(\mathcal{O}_U)$.
Hence we reduce to the case that $X$ is affine, in particular a
quasi-compact and quasi-separated scheme. Via
Lemma \ref{lemma-derived-quasi-coherent-small-etale-site} and
\ref{lemma-descend-perfect}
we reduce to the case of schemes, i.e., to
Derived Categories of Schemes, Proposition
\ref{perfect-proposition-compact-is-perfect}.
\end{proof}

\end{document}
